\documentclass[11pt]{article}
\usepackage[a4paper,margin=25mm]{geometry}
\usepackage{amsmath,amssymb,amsthm,hyperref}
\newtheorem{theorem}{Theorem}
\title{Negative real zeros in the Ma--Wang alternating-run polynomials}
\author{Principia Math In Progress}
\date{22 September 2026}
\begin{document}
\maketitle
\noindent\textbf{Task 1 (MathDB 330839).}
Ma and Wang's published Conjecture 4.1 \cite{MW} asserts that
$T_n(x)/x$ has only imaginary zeros and a coordinate-separation property
for every $n\ge2$. Here ``imaginary'' is read as nonreal, as indicated by
the source's use of general conjugate pairs with nonzero real parts.
The assertion already fails at $n=7$; the failure persists for all larger
indices.

\begin{theorem}
For every integer $n\ge7$, $T_n(x)/x$ has at least two distinct real zeros
in $(-1,0)$ and at least two distinct real zeros in $(-\infty,-1)$.
In particular, the full conjunction in Ma--Wang's Conjecture 4.1 is false.
\end{theorem}
\begin{proof}
The defining recurrence \cite[Eq.~(1.1)]{MW} is
\[
 T_{n+1}(x)=(2nx+1)xT_n(x)+x(1-x^2)T_n'(x),\qquad T_1(x)=x.
\]
Writing $P_n(x)=T_n(x)/x$, we obtain
\begin{equation}\label{eq:rec}
 P_{n+1}(x)=\bigl(1+x+(2n-1)x^2\bigr)P_n(x)
               +x(1-x^2)P_n'(x),\qquad P_1=1.
\end{equation}
In particular,
\begin{equation}\label{eq:end}
 P_n(0)=1,\qquad P_n(-1)=(2n-3)!!>0\quad(n\ge2).
\end{equation}
Six applications of \eqref{eq:rec} give
\begin{align*}
P_7(x)={}&1+63x+1023x^2+4641x^3+16479x^4+26481x^5\\
 &+37759x^6+26481x^7+16479x^8+4641x^9\\
 &+1023x^{10}+63x^{11}+x^{12}.
\end{align*}
Exact substitution yields
\begin{equation}\label{eq:sign}
 P_7(-1/20)=-\frac{316770843280059}{4096000000000000}<0.
\end{equation}
Together with \eqref{eq:end}, the intermediate value theorem already
gives two distinct real zeros in $(-1,0)$.

For persistence, suppose $P_n$ takes a negative value on $(-1,0)$.
By \eqref{eq:end}, its minimum on $[-1,0]$ is attained at an interior
point $a$, with $P_n(a)<0$ and $P_n'(a)=0$. Equation \eqref{eq:rec} gives
\[
 P_{n+1}(a)=\bigl(1+a+(2n-1)a^2\bigr)P_n(a)<0,
\]
because $-1<a<0$ and $n\ge1$. Its endpoint values are again positive,
so it has two distinct zeros in $(-1,0)$. Induction starting from
\eqref{eq:sign} proves this for every $n\ge7$.

The source's symmetry assertion \cite[Theorem 2.2]{MW} gives
\[
 P_n(x)=x^{2n-2}P_n(1/x).
\]
It also follows immediately by induction from the coefficient recurrence
in \cite[Eq.~(1.2)]{MW}. Each of the two zeros in $(-1,0)$ therefore has
a distinct reciprocal zero in $(-\infty,-1)$. The four zeros are distinct
and nonzero, completing the proof.
\end{proof}

\noindent\textbf{Scope and prior work.}
This refutation concerns the nonreal-zero assertion in the original
conjunction. It is compatible with Hurwitz stability, which permits
negative real zeros. Ding and Zhu \cite[Conjecture 2.13 and Proposition
2.16]{DZ} prove Hurwitz stability of the same sequence, while citing
Ma--Wang's Conjecture 4.1 for that different statement. Their result is
credited and is not contradicted here. Their later journal version,
\emph{Advances in Applied Mathematics} 152 (2024), 102591, has not been
read in full; its precise discussion of the original conjecture remains
a priority-audit gap. This note therefore receives \textbf{zero new
campaign credit pending that check}. Both complete versions of the
original Ma--Wang paper were read.

The accompanying exact verifier derives the displayed coefficients by
the recurrence and independently enumerates all $135135$ dual Stirling
permutations of order seven. Rational sign evaluations and Sturm counts
provide additional checks. The persistence argument is analytic and
covers infinitely many indices. No independent peer review or formal
verification is claimed.

\medskip
\noindent\textbf{Proposed Tasks 2--5 (not additional solutions).}
\begin{enumerate}
\setcounter{enumi}{1}
\item Determine the exact number of negative real zeros of $P_n$ for every
$n\ge7$, including multiplicities.
\item Determine whether the proportion of real zeros of $P_n$ has a limit
as $n\to\infty$, and compute it if it exists.
\item Determine the indices for which the real and imaginary coordinates
of the nonreal conjugate pairs of consecutive polynomials interlace.
\item Obtain asymptotic formulas for the negative real zeros closest to
zero and their reciprocal partners.
\end{enumerate}

\begin{thebibliography}{2}
\bibitem{MW} S.-M. Ma and H.-N. Wang,
\emph{Enumeration of a dual set of Stirling permutations by their
alternating runs}, Bulletin of the Australian Mathematical Society
\textbf{94} (2016), 177--186.
\href{https://doi.org/10.1017/S0004972716000186}{doi:10.1017/S0004972716000186}.
Preprint: \href{https://arxiv.org/abs/1506.08716}{arXiv:1506.08716v2},
30 June 2015; its corresponding assertion is Conjecture 7.
\bibitem{DZ} M.-J. Ding and B.-X. Zhu,
\emph{Stability of combinatorial polynomials and its applications},
\href{https://arxiv.org/abs/2106.12176}{arXiv:2106.12176v2}, 24 June 2021.
Related journal version:
\href{https://doi.org/10.1016/j.aam.2023.102591}{doi:10.1016/j.aam.2023.102591}.
\end{thebibliography}
\end{document}
